\subsection*{Chapter 2: Vectors, Bases, and Coordinate Systems}

A coordinate pair locates a point, but it does not yet express how one point is
related to another. The second chapter begins when the question changes from
``Where are we?'' to ``What change carries us from one position to another?'' The
vector is introduced as the answer to that need.

The chapter is meant to show that vector algebra is not a collection of formal
component rules. Addition, scalar multiplication, dot products, projections, and
cross products all answer geometric questions. The algebraic operations are
supposed to feel forced by what displacements must do.

\subsubsection*{Vectors from coordinate displacements}

The first section begins with the difference between a position and a displacement.
Two arrows drawn in different places may represent the same vector because they
record the same change of coordinates. Equality of vectors therefore ignores the
starting point and preserves the displacement itself.

Vector addition is motivated by successive motion: two journeys performed one
after another can be replaced by one net displacement. Scalar multiplication
arises from repeating, shortening, stretching, or reversing a displacement.
Linear combinations combine these two operations and answer a broader question:
what displacements can be built from the directions already available? One
nonzero direction generates a line of possible displacements; two nonparallel
directions generate every displacement in the plane. This idea later becomes the
foundation of bases, lines, and planes.

Length and inclination reconnect coordinate changes with Euclidean geometry. The
dot product is developed as a way of comparing directions, recovering angle, and
detecting perpendicularity. Its algebraic properties are made explicit so that
later expansions are justified rather than performed by habit. Projection then
answers a practical demand: how much of one vector lies in a chosen direction?
The resulting orthogonal decomposition separates a vector into parallel and
perpendicular parts and leads naturally to the Pythagorean relation, the
Cauchy--Schwarz inequality, and the triangle inequality.

The move to three dimensions is intended to show both continuity and novelty.
Addition, scaling, length, and the dot product extend by the same reasoning. The
cross product appears because three-dimensional space permits a new demand: from
two vectors, construct a vector perpendicular to both. Its coordinate formula is
derived from the perpendicularity conditions rather than presented as magic. At
this stage the student should understand two central facts:
\[
\mathbf u\times\mathbf v=\mathbf0
\quad\text{detects parallel directions,}
\]
and
\[
(\mathbf u\times\mathbf v)\cdot\mathbf u
=
(\mathbf u\times\mathbf v)\cdot\mathbf v=0.
\]
The result is therefore perpendicular to both inputs and to every linear
combination of them. The interpretation of those combinations as a plane is
postponed until Chapter 3.

By the end of this section, the student should be able to interpret vector
operations geometrically, move between pictures and component calculations,
construct and analyse linear combinations, compute lengths and angles, test
parallelism and perpendicularity, derive and use projections, and explain what
each product is designed to measure or produce.

\subsubsection*{Changing the basis: new coordinates for the same vector}

Once vectors can be built from two independent directions, a new question becomes
possible: must the horizontal and vertical directions be used to describe every
vector? A basis is a chosen pair of nonparallel directions, and the same geometric
vector receives different coordinate pairs in different bases.

This section deepens the distinction between an object and its coordinates. The
vector remains fixed while the descriptive grid changes. On a rectangular grid
its coordinates are read through a rectangle; on an oblique grid they are read
through a parallelogram. In a non-orthogonal basis, coordinates are not ordinary
perpendicular shadows. They are signed coefficients in
\[
\mathbf u=c_1\mathbf b_1+c_2\mathbf b_2.
\]
The coefficients are found from a system of two linear equations. Matrix notation
is deliberately postponed so that the student first understands the underlying
algebraic problem before later notation compresses it.

The student should learn to determine whether two vectors form a basis, express a
vector in that basis, solve explicitly for the new coordinates, interpret
negative coefficients as motion opposite to a basis direction, and convert
between descriptions without confusing the vector with either coordinate pair.
The broader lesson is that a coordinate system is chosen for a purpose; it is not
imposed by the geometric object itself.

\subsubsection*{Beyond a change of basis: polar coordinates}

After linear changes of coordinates, the final section asks whether every
coordinate system must come from two fixed basis vectors. Polar coordinates show
that the answer is no. A point may instead be described by its distance from the
origin and its direction from the origin.

The contrast is central. A change of basis may slant the grid, but its coordinate
curves remain two families of parallel lines. Polar coordinate curves are circles
and rays, and the direction used to interpret the coordinates varies from point
to point. The formulas
\[
x=r\cos\theta,
\qquad
y=r\sin\theta
\]
therefore describe a nonlinear change of coordinates. Polar pairs cannot be added
componentwise to add vectors.

The student should be able to convert between Cartesian and polar descriptions,
handle the convention at the origin and the chosen angle interval, explain why
$\operatorname{atan2}$ is needed to preserve quadrant information, and recognise
when a nonlinear coordinate system simplifies the geometry. The circle becoming
$r=R$ is the first example of a recurring strategy: choose a language adapted to
the symmetry of the object.
